\section{引言}

\subsection{问题背景}

在电信行业中，获取新客户的成本通常为保留老客户的5至7倍，因此客户流失的有效管控直接关系企业的可持续发展。
某电信运营商近一年内客户流失率由18\%攀升至26\%，亟需识别导致客户流失的关键因素并制定针对性的挽留策略。

本研究基于该运营商提供的7043名客户历史数据，涵盖人口统计学特征、账户服务信息及流失状态记录。
通过数据清洗与探索性分析，建立分类预测模型，结合模型可解释性方法识别核心流失因子，
最终根据不同风险层级制定差异化的客户挽留方案，为运营商降低客户流失率提供数据驱动的决策支持。

\subsection{问题重述}

本文主要通过以下三个递进问题展开分析：

\textbf{问题一：流失客户特征画像。}对原始数据进行预处理和探索性分析，通过可视化手段对比流失客户与未流失客户在各特征维度上的差异，
识别潜在的高相关流失因子。

\textbf{问题二：流失预测与关键因子识别。}构建多种分类模型（逻辑回归、随机森林、XGBoost）对客户流失概率进行预测，
以 AUC 和 F1-score 为评价指标进行模型对比与优选。在此基础上运用 SHAP 方法量化各特征的贡献度，确定前五位关键流失因子。

\textbf{问题三：客户分层与差异化挽留策略。}基于预测流失概率将客户划分为低、中、高三个风险层级，针对不同层级设计差异化的挽留方案，
并结合客户生命周期价值（CLTV）估算策略的预期投入产出比。

\subsection{本文工作概述}

本文综合运用了探索性数据分析、机器学习分类建模和 SHAP 可解释性分析等方法，建立了客户流失预测模型与分层挽留框架，主要工作如下：

\begin{enumerate}
    \item \textbf{针对问题一：}对原始数据进行缺失值检查和数据类型修复，剔除 CustomerID、Count、Country 等无用特征，保留30个可用变量。
    通过计算整体流失率（26.54\%）发现类别不平衡特征。对在网时长（Tenure Months）、月消费金额（Monthly Charges）等数值特征绘制分组直
    方图与箱线图，对合同类型（Contract）、互联网服务类型（Internet Service）等分类特征绘制流失率对比柱状图，
    结合相关性热力图，综合识别出在网时长和合同期限为区分流失客户最显著的特征。
    \item \textbf{针对问题二：}在特征工程基础上，分别构建逻辑回归、随机森林和 XGBoost 三种分类模型，以 AUC 和 F1-score 为核心指标
    进行对比评价。对最优模型使用 SHAP 方法进行全局和局部可解释性分析，量化各特征对流失预测的贡献度，提取前五位关键流失因子，并从业务角
    度进行解读。
    \item \textbf{针对问题三：}根据最优模型的预测流失概率，将客户划分为低风险、中风险和高风险三个层级。针对不同层级客户的行为特征设计
    差异化的挽留策略：对高风险客户给予年付合同优惠和专属客服跟进，对低风险客户维持常规关怀。结合客户生命周期价值估算策略的预期投入产出比，
    论证方案的经济可行性。
\end{enumerate}

最终，通过多模型对比和 SHAP 可解释性分析验证了模型结果的可靠性与可解释性，为运营商提供了切实可行的客户流失管控方案。
